% GENERATED FILE. Edit canonical lessons, then run npm run generate:books.
% canonical-source-hash: fnv1a64:12f6e93075f29e50
% canonical-lessons: ES-C531-desaparecido, ES-C531-busqueda, ES-C531-localizar, ES-C531-hallar, ES-C531-hallazgo

\chapter{La búsqueda}
\label{ch:es-la-busqueda}
\hlchaptermodality{531}

\begin{chapteropening}
\textbf{By the end of this chapter:} \emph{I can say that someone went missing, there was a search, and they were found.}

Complete the last lesson of chapter 531: I can say that someone went missing, there was a search, and they were found.
\end{chapteropening}

\section[desaparecido]{desaparecido — missing; disappeared}

\label{lesson:ES-C531-desaparecido}



\begin{quote}
Before the new one: say the Spanish for the attack, then the Spanish for to attack.
\end{quote}

\subsection*{You'll want to know: desaparecido}
\textbf{desaparecido} — \textquotedblleft{}missing; disappeared\textquotedblright{}.

\begin{quote}
El niño estuvo desaparecido dos días.
\end{quote}

\begin{quote}
\emph{The boy was missing for two days.}
\end{quote}

\emph{Un desaparecido} is a missing person.

\begin{cousinweb}
\emph{desaparecido} is built on the participle of \emph{desaparecer} 'to disappear', from \emph{des-} and \emph{aparecer}, which comes from Latin \emph{appārēre} 'to appear' (Lewis \& Short).
\end{cousinweb}

\begin{grammarlens}[title={What you've built}]
One more word for telling, explaining and describing at B1.
\end{grammarlens}

\subsection*{Your turn}
\begin{itemize}
\raggedright
  \item \emph{Say it:} \emph{desaparecido}
  \item \emph{Say it:} \emph{desaparecido}, once more
  \item \emph{Say it:} a sentence of your own with \emph{desaparecido}
  \item \emph{Recall:} say the Spanish for the attack, then the Spanish for to attack, then say \emph{desaparecido} again
\end{itemize}

\begin{tcolorbox}[breakable,colback=teal!4,colframe=teal!35!black,title={Before you move on}]
What does desaparecido mean? (\textquotedblleft{}Missing; disappeared\textquotedblright{}.) Say it once more.
\end{tcolorbox}
\section[la búsqueda]{la búsqueda — the search}

\label{lesson:ES-C531-busqueda}



\begin{quote}
Before the new one: say the Spanish for to attack, then the Spanish for missing.
\end{quote}

\subsection*{You'll want to know: la búsqueda}
\textbf{la búsqueda} — \textquotedblleft{}the search\textquotedblright{}.

\begin{quote}
La búsqueda duró toda la noche.
\end{quote}

\begin{quote}
\emph{The search lasted all night.}
\end{quote}

\emph{En busca de} — in search of.

\begin{cousinweb}
\emph{búsqueda} is built on \emph{buscar} 'to look for' (EtymDB). The origin of \emph{buscar} itself is not settled.
\end{cousinweb}

\begin{grammarlens}[title={What you've built}]
One more word for telling, explaining and describing at B1.
\end{grammarlens}

\subsection*{Your turn}
\begin{itemize}
\raggedright
  \item \emph{Say it:} \emph{la búsqueda}
  \item \emph{Say it:} \emph{la búsqueda}, once more
  \item \emph{Say it:} a sentence of your own with \emph{la búsqueda}
  \item \emph{Recall:} say the Spanish for to attack, then the Spanish for missing, then say \emph{la búsqueda} again
\end{itemize}

\begin{tcolorbox}[breakable,colback=teal!4,colframe=teal!35!black,title={Before you move on}]
What does la búsqueda mean? (\textquotedblleft{}The search\textquotedblright{}.) Say it once more.
\end{tcolorbox}
\section[localizar]{localizar — to locate, to find; to reach}

\label{lesson:ES-C531-localizar}



\begin{quote}
Before the new one: say the Spanish for missing, then the Spanish for the search.
\end{quote}

\subsection*{You'll want to know: localizar}
\textbf{localizar} — \textquotedblleft{}to locate, to find; to reach\textquotedblright{}.

\begin{quote}
Por fin localizaron el coche.
\end{quote}

\begin{quote}
\emph{At last they located the car.}
\end{quote}

It also means to reach someone by phone: \emph{No te pude localizar.}

\begin{cousinweb}
\emph{localizar} is built on \emph{local}, from Latin \emph{locālis} 'of a place' (Lewis \& Short), from \emph{locus} 'a place', with the ending \emph{-izar}.
\end{cousinweb}

\begin{grammarlens}[title={What you've built}]
One more word for telling, explaining and describing at B1.
\end{grammarlens}

\subsection*{Your turn}
\begin{itemize}
\raggedright
  \item \emph{Say it:} \emph{localizar}
  \item \emph{Say it:} \emph{localizar}, once more
  \item \emph{Say it:} a sentence of your own with \emph{localizar}
  \item \emph{Recall:} say the Spanish for missing, then the Spanish for the search, then say \emph{localizar} again
\end{itemize}

\begin{tcolorbox}[breakable,colback=teal!4,colframe=teal!35!black,title={Before you move on}]
What does localizar mean? (\textquotedblleft{}To locate, to find; to reach\textquotedblright{}.) Say it once more.
\end{tcolorbox}
\section[hallar]{hallar — to find}

\label{lesson:ES-C531-hallar}



\begin{quote}
Before the new one: say the Spanish for the search, then the Spanish for to locate.
\end{quote}

\subsection*{You'll want to know: hallar}
\textbf{hallar} — \textquotedblleft{}to find\textquotedblright{}.

\begin{quote}
Hallaron al niño en el bosque.
\end{quote}

\begin{quote}
\emph{They found the boy in the woods.}
\end{quote}

More formal than \emph{encontrar}; common in the news.

\begin{cousinweb}
\emph{hallar} is inherited from Latin \emph{afflāre} 'to blow, to breathe on' (EtymDB), through old Spanish \emph{fallar}, a hunting word: the dog sniffs and finds.
\end{cousinweb}

\begin{grammarlens}[title={What you've built}]
One more word for telling, explaining and describing at B1.
\end{grammarlens}

\subsection*{Your turn}
\begin{itemize}
\raggedright
  \item \emph{Say it:} \emph{hallar}
  \item \emph{Say it:} \emph{hallar}, once more
  \item \emph{Say it:} a sentence of your own with \emph{hallar}
  \item \emph{Recall:} say the Spanish for the search, then the Spanish for to locate, then say \emph{hallar} again
\end{itemize}

\begin{tcolorbox}[breakable,colback=teal!4,colframe=teal!35!black,title={Before you move on}]
What does hallar mean? (\textquotedblleft{}To find\textquotedblright{}.) Say it once more.
\end{tcolorbox}
\section[el hallazgo]{el hallazgo — the find, the discovery}

\label{lesson:ES-C531-hallazgo}



\begin{quote}
Before the new one: say the Spanish for to locate, then the Spanish for to find.
\end{quote}

\subsection*{You'll want to know: el hallazgo}
\textbf{el hallazgo} — \textquotedblleft{}the find, the discovery\textquotedblright{}.

\begin{quote}
Fue un hallazgo muy importante.
\end{quote}

\begin{quote}
\emph{It was a very important find.}
\end{quote}

\begin{cousinweb}
\emph{hallazgo} is built on \emph{hallar} 'to find', which is inherited from Latin \emph{afflāre} 'to breathe on' (Lewis \& Short).
\end{cousinweb}

\begin{grammarlens}[title={What you've built}]
One more word for telling, explaining and describing at B1.
\end{grammarlens}

\subsection*{Your turn}
\begin{itemize}
\raggedright
  \item \emph{Say it:} \emph{el hallazgo}
  \item \emph{Say it:} \emph{el hallazgo}, once more
  \item \emph{Say it:} a sentence of your own with \emph{el hallazgo}
  \item \emph{Recall:} say the Spanish for to locate, then the Spanish for to find, then say \emph{el hallazgo} again
\end{itemize}

\begin{tcolorbox}[breakable,colback=teal!4,colframe=teal!35!black,title={Before you move on}]
What does el hallazgo mean? (\textquotedblleft{}The find, the discovery\textquotedblright{}.) Say it once more.
\end{tcolorbox}
